\begin{table}[htb!!] % The placement specifier can be [htbp]
\centering
\scalebox{0.85}{
\begin{tcolorbox}[title= Example: Claim from \quantempplus{}] \textbf{Claim}:\texttt{Prime Minister Narendra Modi breached the \textcolor{blue}{election protocol} by addressing a rally in Howrah on \textcolor{orange}{April 6.}}


\textit{\textbf{Ideal Retrieval}}:
\begin{itemize}
    \item \ldots From cooch behar \ldots Howrah \ldots election rally on \textcolor{orange}{April 6}, pm Modi will address\ldots
    \item \ldots Election Silence Period\ldots India (\textcolor{blue}{48 to 24 hours} in advance of polling day and on polling day)\ldots 
    \item \ldots General elections were held in India in seven phases from \textcolor{brown}{11 April to 19 May 2019}\ldots
\end{itemize}
\textit{\textbf{Search Process}}:The given numerical claim requires multiple aspects to be addressed. Firstly, evidence proving that PM Modi did indeed rally on April 6 needs to be fetched. Then, the information on the election silence period and the polling schedule are to be retrieved. Given this information, the downstream verifier can then reason about the numerical aspects to form the verdict.


\end{tcolorbox}
}

\captionof{figure}{Example numerical claim from \quantempplus{}, requiring to recognize and fetch explicit and implicit evidence.}\label{intro:example} 
\end{table}%